\documentclass[a4paper,12pt]{article}

\usepackage[T2A]{fontenc}
\usepackage[utf8]{inputenc}
\usepackage[russian]{babel}
\usepackage{amsmath,amssymb,amsthm,mathtools}
\usepackage{helvet}
\renewcommand{\familydefault}{\sfdefault}
\usepackage{icomma}
\usepackage[a4paper,margin=2.2cm,headheight=15pt]{geometry}
\usepackage{microtype}
\usepackage{tikz}
\usetikzlibrary{calc}
\usepackage{tkz-euclide}
\usepackage{pgfplots}
\pgfplotsset{compat=1.18}
\usepackage{tabularx,booktabs,longtable}
\usepackage{enumitem}
\usepackage{hyperref}
\usepackage{parskip}
\usepackage{fancyhdr}
\usepackage{setspace}
\usepackage{xcolor}

\definecolor{accent}{HTML}{008080}
\definecolor{darkaccent}{HTML}{004D4D}
\definecolor{textgray}{HTML}{333333}
\definecolor{softgray}{HTML}{F1F4F4}

\hypersetup{
  colorlinks=true,
  linkcolor=darkaccent,
  urlcolor=darkaccent,
  pdftitle={Чек-лист: прикладные модели, степени и линейная функция},
  pdfauthor={Лёвин Артём Александрович}
}

\setstretch{1.5}
\setlength{\parindent}{0pt}

\setlist[itemize]{
  leftmargin=1.4em,
  itemsep=0.25em,
  topsep=0.35em
}

\setlist[enumerate]{
  leftmargin=1.7em,
  itemsep=0.35em,
  topsep=0.35em
}

\renewcommand{\arraystretch}{1.28}

\pagestyle{fancy}
\fancyhf{}
\fancyhead[L]{\small\textcolor{textgray}{Ксения Клыкова · 10 класс}}
\fancyhead[R]{\small\textcolor{textgray}{Повторение}}
\fancyfoot[C]{\small\thepage}
\renewcommand{\headrulewidth}{0.4pt}
\renewcommand{\headrule}{%
  \hbox to\headwidth{%
    \color{accent!55}\leaders\hrule height \headrulewidth\hfill
  }%
}

\newcommand{\sectionline}{%
  \par
  \vspace{-0.25em}
  \textcolor{accent!65}{\rule{\linewidth}{0.6pt}}
  \vspace{0.35em}
}

\newcommand{\term}[1]{\textbf{\textcolor{darkaccent}{#1}}}

\begin{document}

\begin{titlepage}
  \thispagestyle{empty}
  \centering
  \vspace*{2.5cm}

  {\Large\textcolor{darkaccent}{ЧЕК-ЛИСТ}\par}

  \vspace{0.7cm}

  {\huge\bfseries Прикладные модели, степени\par}

  \vspace{0.15cm}

  {\huge\bfseries и линейная функция\par}

  \vspace{0.8cm}

  {\large Повторение ключевых методов решения задач\par}

  \vfill

  {\large Лёвин Артём Александрович\par}

  \vspace{0.15cm}

  {\large эксклюзивно для Ксении Клыковой\par}

  \vspace{0.45cm}

  {\large \today\par}

  \begin{tikzpicture}[remember picture,overlay]
    \draw[
      accent!48,
      line width=1.2pt
    ]
      ($(current page.south west)+(1.5cm,1.45cm)$)
      .. controls
      ($(current page.south west)+(5.0cm,2.10cm)$)
      and
      ($(current page.south east)+(-5.0cm,0.85cm)$)
      ..
      ($(current page.south east)+(-1.5cm,1.50cm)$);

    \draw[
      darkaccent!30,
      line width=0.55pt
    ]
      ($(current page.south west)+(1.5cm,1.18cm)$)
      .. controls
      ($(current page.south west)+(5.5cm,1.75cm)$)
      and
      ($(current page.south east)+(-5.5cm,0.65cm)$)
      ..
      ($(current page.south east)+(-1.5cm,1.20cm)$);
  \end{tikzpicture}
\end{titlepage}

\section*{1. Главная стратегия}

\sectionline

В задачах прикладного характера сначала переводят текст на язык
формул, затем решают полученную математическую модель и проверяют
смысл ответа.

\begin{enumerate}
  \item Выпишите величины и кратко подпишите их смысл.

  \item Запишите основную формулу или зависимость.

  \item Выразите искомую величину через известные данные.

  \item Зафиксируйте ограничения: цена, время, объём и температура
  обычно принимают допустимые по смыслу значения.

  \item Выполните вычисления, контролируя знаки, степени и единицы
  измерения.

  \item Проверьте результат подстановкой и выберите ответ,
  подходящий условию.
\end{enumerate}

\begin{table}[ht]
  \centering
  \caption{Обозначения, использованные в чек-листе}

  \begin{tabularx}{\linewidth}{@{}lX@{}}
    \toprule
    \textbf{Обозначение} & \textbf{Смысл} \\
    \midrule
    $p$   & цена единицы продукции \\
    $Q$   & объём производства, продаж или спроса \\
    $v$   & переменные затраты на единицу продукции \\
    $F$   & постоянные расходы \\
    $\Pi$ & прибыль \\
    $R$   & выручка \\
    $P$   & мощность в физической формуле \\
    $T$   & абсолютная температура \\
    \bottomrule
  \end{tabularx}
\end{table}

\section*{2. Прибыль: формула и алгоритм}

\sectionline

\term{Базовая модель прибыли:}
\[
  \Pi=Q(p-v)-F.
\]

Величина $p-v$ показывает прибыль с одной единицы продукции до учёта
постоянных расходов.

\textbf{Пример.}
Цена равна $500$, переменные затраты равны $300$, постоянные расходы
составляют $700\,000$, требуемая прибыль равна $300\,000$.
Найдём объём $Q$:
\[
  300\,000=Q(500-300)-700\,000,
\]
\[
  200Q=1\,000\,000,
  \qquad
  Q=5\,000.
\]

\term{Контроль смысла:}
объём производства должен быть положительным; при целочисленном
выпуске итоговый ответ задают целым числом единиц.

\section*{3. Спрос и выручка}

\sectionline

Если спрос зависит от цены, часто используют линейную модель
\[
  Q=a-bp,
  \qquad
  a>0,
  \qquad
  b>0.
\]

Выручка определяется формулой
\[
  R=pQ=p(a-bp).
\]

Слова \textit{«не менее»} переводятся знаком $\ge$, слова
\textit{«не более»} --- знаком $\le$.
Равенство помогает найти граничные значения, после чего требуется
решить всё неравенство.

\textbf{Пример.}
Пусть
\[
  Q=100-10p,
  \qquad
  R\ge 240.
\]

Тогда
\[
  p(100-10p)\ge 240,
\]
\[
  -10p^2+100p-240\ge 0,
\]
\[
  p^2-10p+24\le 0.
\]

При делении на отрицательное число знак неравенства меняется.
Получаем
\[
  (p-4)(p-6)\le 0,
  \qquad
  p\in[4;6].
\]

Наибольшая допустимая цена равна
\[
  \boxed{p=6}.
\]

Дополнительно проверяются экономические ограничения
\[
  p\ge 0,
  \qquad
  Q\ge 0.
\]

\section*{4. Неявные данные и физический смысл}

\sectionline

Часть данных может содержаться в формулировке задачи.

\begin{itemize}
  \item Бак опустел: высота воды равна $h=0$.

  \item Тело остановилось: скорость равна $v=0$.

  \item Производство окупилось: прибыль равна $\Pi=0$.

  \item Объём, время, масса, цена и абсолютная температура выбираются
  с учётом физического или экономического смысла.
\end{itemize}

Если квадратное уравнение даёт несколько корней, каждый корень
проверяют по условию.
Отрицательное время и отрицательный объём обычно исключаются.

\section*{5. Стандартный вид числа и работа со степенями}

\sectionline

\term{Стандартный вид:}
\[
  a\cdot 10^n,
  \qquad
  1\le a<10,
  \qquad
  n\in\mathbb{Z}.
\]

Полезные правила:
\[
  10^m\cdot 10^n=10^{m+n},
  \qquad
  \frac{10^m}{10^n}=10^{m-n},
  \qquad
  \left(10^m\right)^r=10^{mr}.
\]

При переносе десятичной запятой вправо показатель степени уменьшается
на число перемещений:
\[
  5{,}7\cdot 10^{-8}=57\cdot 10^{-9},
  \qquad
  9{,}12\cdot 10^{25}=912\cdot 10^{23}.
\]

\subsection*{Формула Стефана--Больцмана}

\[
  P=\sigma ST^4.
\]

Здесь $P$ --- мощность, $\sigma$ --- постоянная, $S$ --- площадь
поверхности, $T$ --- абсолютная температура.

\textbf{Пример.}
Дано
\[
  P=9{,}12\cdot 10^{25},
  \qquad
  \sigma=5{,}7\cdot 10^{-8},
  \qquad
  S=\frac{1}{16}\cdot 10^{20}.
\]

Тогда
\[
  T^4=\frac{P}{\sigma S}
  =
  \frac{
    912\cdot 10^{23}\cdot 16
  }{
    57\cdot 10^{-9}\cdot 10^{20}
  }.
\]

Поскольку $912=16\cdot 57$, получаем
\[
  T^4=2^8\cdot 10^{12}.
\]

Абсолютная температура положительна, поэтому
\[
  T=\sqrt[4]{2^8\cdot 10^{12}}
  =2^2\cdot 10^3
  =4\,000.
\]

\term{Практический приём:}
большие коэффициенты раскладывайте на множители до извлечения корня.
Такой способ сокращает объём арифметики.

\section*{6. Линейная функция и метод узловых точек}

\sectionline

Линейная функция имеет вид
\[
  f(x)=kx+b.
\]

Для нахождения двух коэффициентов $k$ и $b$ нужны две точки графика
с различными абсциссами.

Пусть известны точки
\[
  A(x_1,y_1),
  \qquad
  B(x_2,y_2),
  \qquad
  x_1\ne x_2.
\]

Тогда можно использовать систему
\[
  \begin{cases}
    y_1=kx_1+b,\\
    y_2=kx_2+b
  \end{cases}
\]
или формулы
\[
  k=\frac{y_2-y_1}{x_2-x_1},
  \qquad
  b=y_1-kx_1.
\]

\textbf{Пример.}
График проходит через точки
\[
  A(3;4),
  \qquad
  B(-1;-3).
\]

Получаем систему
\[
  \begin{cases}
    3k+b=4,\\
    -k+b=-3.
  \end{cases}
\]

Отсюда
\[
  k=\frac{7}{4},
  \qquad
  b=-\frac{5}{4},
\]
следовательно,
\[
  f(x)=\frac{7}{4}x-\frac{5}{4}.
\]

Тогда
\[
  f(-5)
  =
  \frac{7}{4}\cdot(-5)-\frac{5}{4}
  =
  -\frac{40}{4}
  =
  -10.
\]

\begin{figure}[ht]
  \centering

  \begin{tikzpicture}
    \begin{axis}[
      width=0.78\linewidth,
      height=0.50\linewidth,
      axis lines=middle,
      unit vector ratio=1 1,
      xmin=-5.5,
      xmax=4.5,
      ymin=-11,
      ymax=7,
      samples=220,
      clip=false,
      grid=both,
      minor tick num=1,
      xlabel={$x$},
      ylabel={$y$},
      tick label style={font=\small},
      label style={font=\small},
      grid style={draw=softgray},
      axis line style={draw=textgray}
    ]
      \addplot[
        domain=-5.5:4.5,
        accent,
        thick
      ]
      {7*x/4-5/4};

      \addplot[
        only marks,
        mark=*,
        mark size=2.4pt,
        darkaccent
      ]
      coordinates {
        (3,4)
        (-1,-3)
      };

      \node[
        anchor=south west,
        text=darkaccent
      ]
      at (axis cs:3,4)
      {$A(3;4)$};

      \node[
        anchor=north east,
        text=darkaccent
      ]
      at (axis cs:-1,-3)
      {$B(-1;-3)$};
    \end{axis}
  \end{tikzpicture}

  \caption{Линейная функция, восстановленная по двум узловым точкам}
\end{figure}

\section*{7. Типичные ошибки}

\sectionline

\begin{enumerate}
  \item Подмена неравенства одним равенством. После нахождения границ
  требуется определить весь промежуток решений.

  \item Пропуск ограничений
  \[
    p\ge 0,
    \qquad
    Q\ge 0,
    \qquad
    t\ge 0,
    \qquad
    T>0.
  \]

  \item Подстановка выражения без скобок, особенно при наличии знака
  минус.

  \item Сохранение прежнего знака неравенства после деления на
  отрицательное число.

  \item Ошибка в показателях при делении степеней: из верхнего
  показателя вычитают нижний.

  \item Преждевременное округление. Точные дроби сохраняют до
  последнего шага.

  \item Сложение дробей с разными знаменателями без приведения к
  общему знаменателю.

  \item Потеря знака при умножении двух отрицательных чисел.

  \item Перестановка координат точки: сначала записывают $x$, затем
  $y=f(x)$.

  \item Выбор корня без проверки по смыслу задачи.
\end{enumerate}

\section*{8. Итоговый чек-лист перед записью ответа}

\sectionline

\begin{itemize}
  \item Все величины подписаны и различимы по обозначениям.

  \item Текст задачи переведён в формулу, уравнение или неравенство.

  \item Искомая величина осталась единственной неизвестной.

  \item Скобки при подстановке расставлены.

  \item Знаки при переносе и делении проверены.

  \item Степени приведены к единому основанию, коэффициенты сокращены.

  \item Для дробей найден общий знаменатель.

  \item Все корни проверены по ограничениям и смыслу.

  \item Ответ записан в требуемой форме: число, промежуток, цена,
  время или объём.
\end{itemize}

\clearpage

\section*{Тренировочный блок}

\sectionline

\subsection*{Средний уровень}

\begin{enumerate}[label=\textbf{\arabic*.}]
  \item Цена единицы продукции равна $750$ рублей, переменные затраты
  равны $430$ рублей, постоянные расходы составляют $960\,000$ рублей.
  Какой объём выпуска обеспечит прибыль $640\,000$ рублей?

  \item Спрос описывается формулой
  \[
    Q=120-8p.
  \]
  Выручка должна составить не менее $400$.
  Найдите наибольшую допустимую цену $p$.

  \item Высота воды в баке задаётся формулой
  \[
    h(t)=-t^2+12t+13,
  \]
  где $t\ge 0$.
  Через сколько единиц времени бак опустеет?
\end{enumerate}

\subsection*{Уровень выше среднего}

\begin{enumerate}[label=\textbf{\arabic*.},resume]
  \item В формуле
  \[
    P=\sigma ST^4
  \]
  даны
  \[
    \sigma=5{,}7\cdot 10^{-8},
    \qquad
    S=10^{14},
    \qquad
    P=9{,}12\cdot 10^{19}.
  \]
  Найдите $T$.

  \item Линейная функция проходит через точки
  $A(2;5)$ и $B(-2;-3)$.
  Найдите $f(6)$.

  \item Линейная функция проходит через точки
  $A(4;1)$ и $B(-4;-2)$.
  Найдите $f(-10)$.

  \item Цена единицы продукции равна $900$ рублей, переменные затраты
  равны $540$ рублей, постоянные расходы составляют
  $1\,260\,000$ рублей.
  Найдите наименьший целый объём выпуска, при котором прибыль будет
  не менее $900\,000$ рублей.

  \item Спрос описывается формулой
  \[
    Q=180-6p.
  \]
  Выручка должна быть не менее $1\,200$.
  Найдите наибольшую допустимую цену.

  \item Известно, что
  \[
    f(3)=7,
    \qquad
    f(-5)=-9,
  \]
  причём
  \[
    f(x)=kx+b.
  \]
  Найдите значение $x$, при котором $f(x)=25$.
\end{enumerate}

\subsection*{Сложный уровень}

\begin{enumerate}[label=\textbf{\arabic*.},resume]
  \item Зависимость спроса $Q$ от цены $p$ линейна.
  Известны две точки:
  \[
    (p;Q)=(5;150),
    \qquad
    (p;Q)=(20;60).
  \]

  Переменные затраты составляют $5$ денежных единиц на единицу
  продукции, постоянные расходы равны $300$.
  Прибыль определяется формулой
  \[
    \Pi=(p-5)Q-300.
  \]

  Найдите наибольшую целую цену $p$, при которой прибыль составляет
  не менее $300$, и укажите соответствующий объём спроса.
\end{enumerate}

\clearpage

\section*{Ответы к упражнениям}

\sectionline

\begin{table}[ht]
  \centering
  \caption{Краткие ответы}

  \begin{tabularx}{\linewidth}{@{}cX@{}}
    \toprule
    \textbf{№} & \textbf{Ответ} \\
    \midrule
    1  & $Q=5\,000$ \\
    2  & $p=10$ \\
    3  & $t=13$ \\
    4  & $T=2\,000$ \\
    5  & $f(6)=13$ \\
    6  & $f(-10)=-\dfrac{17}{4}$ \\
    7  & $Q=6\,000$ \\
    8  & $p=20$ \\
    9  & $x=12$ \\
    10 & $p=25$, соответствующий спрос $Q=30$ \\
    \bottomrule
  \end{tabularx}
\end{table}

\vfill

\begin{center}
  \textcolor{darkaccent}{
    \small
    Проверяйте модель, вычисления и смысл результата на каждом этапе.
  }
\end{center}

\end{document}